Normal-only video anomaly detectors are often compared by test-set ranking metrics, although deployment also requires a threshold that survives a change of camera or dataset. We test whether a low-capacity scene-residual representation can improve this transfer. The tested Scene-Residual Normality (SRN) module predicts a scene component from frozen DINOv2 features, subtracts it, and optionally restores a controlled context vector. We evaluate official UCSD Ped2 and CUHK Avenue ground truth with whole-video splits, source-normal threshold selection, a separate target-normal calibration track, three normality scorers, and mechanism probes. Cross-dataset threshold tests use raw frozen-feature scorers; SRN is evaluated only in a joint two-seen-domain mechanism diagnostic. The negative result is clear. Full SRN reaches frame AUROC $0.6677$, below raw Gaussian scoring at $0.6885$ and only $0.0024$ above its matched raw-prototype head. Dataset/camera identity remains perfectly decodable from the SRN residual. More seriously, every tested source-normal threshold yields target-normal false-positive rate $1.0$ in both cross-dataset directions. Oracle test-curve operating points and target-normal calibration are reported separately and do not repair this deployment failure. These experiments falsify the tested low-capacity SRN mechanism on Ped2/Avenue; they do not test genuine unseen-scene episodic leave-one-scene-out generalization. The study shows why representation-removal claims require both an independent identity probe and fixed-threshold evaluation.
